\documentclass[10pt,a4paper]{amsart}
\usepackage[margin=19mm]{geometry}
\usepackage{amsmath,amssymb,mathtools}
\usepackage[scr=rsfs,cal=euler]{mathalfa}
\usepackage{xfrac}
\usepackage[unicode,hidelinks,bookmarks=false]{hyperref}
\newcommand{\ie}{i.e.\ }
\newcommand{\cf}{cf.\ }
\newcommand{\resp}{respectively\ }
\newcommand{\Subsection}{\subsection}
\providecommand{\nobreakdash}{\nobreak\hbox{-}}
\setlength{\emergencystretch}{2em}
\multlinegap0pt
\usepackage[shortlabels]{enumitem}
\usepackage{amssymb}
\usepackage{mathtools}
\usepackage{dsfont}
\usepackage{bm}
\newtheorem{theorem}{Theorem}
\def\C			{\mathds{C}}
\DeclareMathOperator{\diag}{diag}
\def\im	 					{\mathrm{i}}				% imaginary unit
\title{Hermitian Geometry of Complex Multivectors, Determinants and Orientations}
\author{Andr\'e L. G. Mandolesi}
\date{Source: arXiv:2403.17022v4 (CC BY 4.0)}
\begin{document}
\maketitle

\noindent\emph{Source and licence.} Hermitian Geometry of Complex Multivectors, Determinants and Orientations. Version arXiv:2403.17022v4 (CC BY 4.0), distributed by arXiv under the Creative Commons Attribution 4.0 International licence, http://creativecommons.org/licenses/by/4.0/, as recorded in the arXiv licence metadata for this exact version and shown on the abstract page https://arxiv.org/abs/2403.17022v4. Full licence notice: \texttt{licenses/arxiv-2403.17022-CC-BY-4.0.txt}.\par\smallskip

\noindent\textit{Editorial scope.} Counted unit: the article's theorem pr:decomposition (source0.tex lines 1005--1016). The article's complete original proof of it is delivered with it (source0.tex lines 1017--1054, one \texttt{proof} environment of 38 lines). No mathematics is written, repaired or re-derived: the statement and the proof are the article's own lines, in the article's own notation, and no retained line is substituted or re-flowed.  No further source-local context is needed: every cross-reference of every retained line resolves inside the two delivered spans.  Nothing was omitted from any retained span, so the package carries no omission record.  No retained line contains a citation, so the wrapper carries no bibliography and no bibliography entry.  No retained line contains a figure environment, an image-inclusion command or drawing code, so no image is delivered and the package declares no asset dependency.  Editorial formatting only, in addition: an AMS \texttt{amsart} preamble that restates the article's own declarations and macros used by the retained lines, namely the environments \texttt{theorem} on separate counters, together with the standard packages those lines need.  The retained environments print in this document as Theorem 1 in order of appearance, whereas the article prints the same items with the numbering of its own full document.  The complete native source of the article as distributed in the arXiv e-print for this exact version is bundled unchanged as \texttt{sources/156625/source0.tex}.
\begin{theorem}\label{pr:decomposition}
	Let $\beta = (v_1,\ldots,v_p)$ be a basis of $V \subset \C^n$ in which $T \in GL(V)$ has Jordan normal form, with eigenvalues $\lambda_1,\ldots,\lambda_p$ in its diagonal. 
	Then:
	\begin{enumerate}[label = (\roman*)]
		\item $T$ is a composition of complex shears, line scalings and phase rotations by $\varphi_j = \arg(\lambda_j)$ of the $\C v_j$'s,
		and $\arg(\det T) = \sum_{j=1}^p \varphi_j \pmod{2\pi}$. \label{it:T phases}
		
		\item $T$ is a composition of complex shears, line scalings and a single phase rotation by $\varphi = \arg(\det T)$. \label{it:T single phase}
		
		\item $T \in GL^+(V) \Leftrightarrow T$ is a composition of complex shears and line scalings. \label{it:GL+}
	\end{enumerate}
\end{theorem}

\begin{proof}
	\ref{it:T phases} 
	Each $\lambda_j$ rotates $\C v_j$ by $\varphi_j$ and scales it by $|\lambda_j|$,
	any superdiagonal $1$ causes a complex shear, and $\arg(\det T) = \arg\big(\prod_{j=1}^p \lambda_j\big) = \sum_{j=1}^p \varphi_j \pmod{2\pi}$.
	
	\ref{it:T single phase} The phase rotations in \ref{it:T phases} are produced by $R = \diag(e^{\im \varphi_1},\ldots,e^{\im \varphi_p})$ in the basis $\beta$. If $p=2$ we have, with $a = e^{\im \varphi_1}$,
	\begin{equation*}
		\begin{pmatrix}
			e^{\im \varphi_1} & 0 \\
			0 & e^{\im \varphi_2}
		\end{pmatrix} =
		\begin{pmatrix}
			1 & 0 \\
			0 & e^{\im (\varphi_1+\varphi_2)}
		\end{pmatrix}
		\begin{pmatrix}
			1 & 1-a \\
			0 & 1
		\end{pmatrix}
		\begin{pmatrix}
			1 & 0 \\
			-1 & 1
		\end{pmatrix}
		\begin{pmatrix}
			1 & 1-\bar{a} \\
			0 & 1
		\end{pmatrix}
		\begin{pmatrix}
			1 & 0 \\
			a & 1
		\end{pmatrix},	
	\end{equation*}
	so $R$ consists of complex shears and a single phase rotation of $\C v_2$ by $\varphi_1+\varphi_2$.
	If $p>2$, doing the same 2 phases at a time we can move them down the diagonal to obtain $R = \diag(1,\ldots,1,e^{\im\varphi}) \circ S$, where $\varphi = \sum_{j=1}^p \varphi_j = \arg(\det T)$
	and $S$ is a composition of complex shears.
	
	\ref{it:GL+} Follows from \ref{it:T single phase}.
\end{proof}

\end{document}
